\documentclass[11pt]{article}
% ======================================================================
%  Preambule commun - Sujets Bac Serie C (Physique-Chimie)
%  Style sobre : noir et blanc, sans couleurs, sans filigrane,
%  sans en-tetes ni pieds decoratifs. Figures reconstruites en TikZ.
% ======================================================================
\usepackage[utf8]{inputenc}
\usepackage[T1]{fontenc}
\usepackage{lmodern}
\usepackage{textcomp}
\usepackage{amsmath,amssymb}
\usepackage{enumitem}
\usepackage[a4paper,margin=2.3cm]{geometry}
\usepackage{fancyhdr}
\usepackage{tikz}
\usetikzlibrary{arrows.meta,calc,decorations.pathmorphing,patterns,angles,quotes}

% Symbole degre robuste + justification souple
\DeclareUnicodeCharacter{00B0}{\ensuremath{^\circ}}
\tolerance=2000
\emergencystretch=2em
\frenchspacing

% En-tete / pied minimalistes
% Pied de page (provenance) + entete corrige + encadres pedagogiques
% ======================================================================
%  Fragment commun a tous les styles (sujets et corriges).
%  Inclus depuis chaque dossier serie : % ======================================================================
%  Fragment commun a tous les styles (sujets et corriges).
%  Inclus depuis chaque dossier serie : \input{../../commun}
%  (la compilation se fait toujours depuis le dossier de la serie).
%
%  Style sobre conserve : noir et blanc uniquement, filets fins,
%  pas d'aplats ni de couleurs — lisible en photocopie.
% ======================================================================

% --- Compression maximale des PDF ------------------------------------
% Le public consulte souvent sur reseaux mobiles : chaque Ko compte.
\pdfcompresslevel=9
\pdfobjcompresslevel=3
\pdfminorversion=5

% --- Adresse publique du site ---------------------------------------
% Point unique a mettre a jour si le domaine change (puis recompiler
% tous les PDF : voir sources/README.md).
\newcommand{\bacsiteurl}{annabac.pages.dev}

% --- Pied de page : provenance discrete + numero de page ------------
% Les PDF circulent detaches du site (messageries, photocopies) : cette
% ligne permet de retrouver la source et rappelle la gratuite.
\pagestyle{fancy}
\fancyhf{}
\renewcommand{\headrulewidth}{0pt}
\fancyfoot[L]{\scriptsize\itshape Bac 242 --- annales gratuites du bac congolais --- \bacsiteurl}
\fancyfoot[R]{\thepage}

% --- Entete structure des corriges -----------------------------------
% Usage : \entetecorrige{2020}{C}{Mathématiques}
% Les sujets, reproductions de documents officiels, gardent \entetesujet.
\newcommand{\entetecorrige}[3]{%
  \begin{center}
    {\scshape Baccalaur\'eat --- S\'erie #2}\\[0.35em]
    {\Large\bfseries #3 --- Session #1}\\[0.4em]
    {\itshape Corrig\'e d\'etaill\'e}
  \end{center}
  \vspace{0.2em}
  \hrule height 0.8pt
  \vspace{1.5pt}
  \hrule height 0.4pt
  \vspace{1.2em}}

% --- Encadres pedagogiques (corriges) --------------------------------
% Filet vertical gauche, texte legerement reduit, aucun fond : sobre,
% photocopiable, et tolere les sauts de page (package framed).
% Usage, avec parcimonie (2-3 encadres par corrige, la ou ils eclairent
% une demarche) :
%   \begin{methode} ... \end{methode}   -> « Méthode »
%   \begin{rappel}  ... \end{rappel}    -> « Rappel de cours »
%   \begin{piege}   ... \end{piege}     -> « Piège classique »
\usepackage{framed}
\newenvironment{pedagobarre}{%
  \def\FrameCommand{{\vrule width 1pt}\hspace{0.9em}}%
  \MakeFramed{\advance\hsize-\width\FrameRestore}}%
 {\endMakeFramed}
\newenvironment{blocpedago}[1]{%
  \par\medskip
  \begin{pedagobarre}\small
  \noindent{\bfseries\scshape #1.}\hspace{0.5em}\ignorespaces
}{%
  \end{pedagobarre}\medskip}
\newenvironment{methode}{\begin{blocpedago}{M\'ethode}}{\end{blocpedago}}
\newenvironment{rappel}{\begin{blocpedago}{Rappel de cours}}{\end{blocpedago}}
\newenvironment{piege}{\begin{blocpedago}{Pi\`ege classique}}{\end{blocpedago}}

%  (la compilation se fait toujours depuis le dossier de la serie).
%
%  Style sobre conserve : noir et blanc uniquement, filets fins,
%  pas d'aplats ni de couleurs — lisible en photocopie.
% ======================================================================

% --- Compression maximale des PDF ------------------------------------
% Le public consulte souvent sur reseaux mobiles : chaque Ko compte.
\pdfcompresslevel=9
\pdfobjcompresslevel=3
\pdfminorversion=5

% --- Adresse publique du site ---------------------------------------
% Point unique a mettre a jour si le domaine change (puis recompiler
% tous les PDF : voir sources/README.md).
\newcommand{\bacsiteurl}{annabac.pages.dev}

% --- Pied de page : provenance discrete + numero de page ------------
% Les PDF circulent detaches du site (messageries, photocopies) : cette
% ligne permet de retrouver la source et rappelle la gratuite.
\pagestyle{fancy}
\fancyhf{}
\renewcommand{\headrulewidth}{0pt}
\fancyfoot[L]{\scriptsize\itshape Bac 242 --- annales gratuites du bac congolais --- \bacsiteurl}
\fancyfoot[R]{\thepage}

% --- Entete structure des corriges -----------------------------------
% Usage : \entetecorrige{2020}{C}{Mathématiques}
% Les sujets, reproductions de documents officiels, gardent \entetesujet.
\newcommand{\entetecorrige}[3]{%
  \begin{center}
    {\scshape Baccalaur\'eat --- S\'erie #2}\\[0.35em]
    {\Large\bfseries #3 --- Session #1}\\[0.4em]
    {\itshape Corrig\'e d\'etaill\'e}
  \end{center}
  \vspace{0.2em}
  \hrule height 0.8pt
  \vspace{1.5pt}
  \hrule height 0.4pt
  \vspace{1.2em}}

% --- Encadres pedagogiques (corriges) --------------------------------
% Filet vertical gauche, texte legerement reduit, aucun fond : sobre,
% photocopiable, et tolere les sauts de page (package framed).
% Usage, avec parcimonie (2-3 encadres par corrige, la ou ils eclairent
% une demarche) :
%   \begin{methode} ... \end{methode}   -> « Méthode »
%   \begin{rappel}  ... \end{rappel}    -> « Rappel de cours »
%   \begin{piege}   ... \end{piege}     -> « Piège classique »
\usepackage{framed}
\newenvironment{pedagobarre}{%
  \def\FrameCommand{{\vrule width 1pt}\hspace{0.9em}}%
  \MakeFramed{\advance\hsize-\width\FrameRestore}}%
 {\endMakeFramed}
\newenvironment{blocpedago}[1]{%
  \par\medskip
  \begin{pedagobarre}\small
  \noindent{\bfseries\scshape #1.}\hspace{0.5em}\ignorespaces
}{%
  \end{pedagobarre}\medskip}
\newenvironment{methode}{\begin{blocpedago}{M\'ethode}}{\end{blocpedago}}
\newenvironment{rappel}{\begin{blocpedago}{Rappel de cours}}{\end{blocpedago}}
\newenvironment{piege}{\begin{blocpedago}{Pi\`ege classique}}{\end{blocpedago}}


% Macros maths / physique
\newcommand{\vv}[1]{\overrightarrow{#1}}
\newcommand{\R}{\mathbb{R}}
\newcommand{\N}{\mathbb{N}}
\newcommand{\vi}{\vec{\imath}}
\newcommand{\vj}{\vec{\jmath}}
% Chimie : formules en romain
\newcommand{\ch}[1]{\ensuremath{\mathrm{#1}}}

% Titre de sujet
\newcommand{\entetesujet}[1]{\begin{center}{\Large\bfseries #1}\end{center}\vspace{0.3em}\hrule\vspace{1em}}

% Bandeau de matiere : CHIMIE / PHYSIQUE avec bareme
\newcommand{\matiere}[2]{\par\bigskip\begin{center}\rule{2.2cm}{0.4pt}\quad{\large\bfseries #1}\ \textit{#2}\quad\rule{2.2cm}{0.4pt}\end{center}\nopagebreak\medskip}

% Titres d'exercices et parties
\newcommand{\exo}[2]{\medskip\noindent{\large\bfseries Exercice #1}\hfill\textit{#2}\par\nopagebreak\smallskip}
\newcommand{\partie}[1]{\medskip\noindent{\bfseries Partie #1}\par\nopagebreak\smallskip}
% Titre de question (corriges) : en gras, sur sa propre ligne
\newcommand{\question}[1]{\par\medskip\noindent{\bfseries #1}\par\nopagebreak\smallskip}

\setlist{leftmargin=2em,itemsep=2pt,topsep=3pt}


\begin{document}

\entetecorrige{2016}{C}{Physique-Chimie}

\matiere{CHIMIE}{8 points}
\partie{A : vérification des connaissances}

\question{1. Vrai ou faux}
\textbf{1. a. Vrai.} La saponification résulte de l'action de l'ion hydroxyde \ch{HO^-} sur un ester ; elle conduit à un alcool et à un ion carboxylate (base conjuguée d'un acide carboxylique) :
\[ \underset{\text{ester}}{\ch{R{-}CO{-}O{-}R'}} + \ch{HO^-} \longrightarrow \underset{\text{ion carboxylate}}{\ch{R{-}CO{-}O^-}} + \underset{\text{alcool}}{\ch{R'OH}} \]
La réaction de saponification est \textbf{lente} et \textbf{totale}.

\textbf{1. b. Faux.} Pour un acide faible de concentration $C$ : $\text{pH} = \frac{1}{2}(\text{p}K_a - \log C)$. Après dilution d'un facteur $N > 1$ ($C' = C/N$, $\text{p}K_a$ inchangé) :
\[ \text{pH}' = \tfrac{1}{2}(\text{p}K_a - \log C') = \tfrac{1}{2}(\text{p}K_a - \log C + \log N) = \text{pH} + \tfrac{1}{2}\log N. \]
On en déduit $\text{pH}' > \text{pH}$.

\question{2. Texte à trous}
Une réaction d'oxydoréduction est une réaction de \emph{transfert} d'électrons au cours de laquelle il y a simultanément \emph{oxydation} du réducteur et \emph{réduction} de l'oxydant.

\question{3. Appariement}
$\mathbf{a_1 = b_3}$ : énergie nécessaire pour ioniser un atome d'hydrogène depuis son niveau fondamental (arracher un électron), valant 13,6 eV.

$\mathbf{a_2 = b_2}$ : toute particule de masse $m$ au repos possède une énergie $E_0 = mc^2$ (équivalence masse-énergie).

$\mathbf{a_3 = b_4}$ : la pile Daniell cuivre-zinc est constituée de deux demi-piles reliées par un pont salin (ici nitrate d'ammonium).
\begin{center}
\begin{tikzpicture}[>=Stealth,scale=0.9,every node/.style={font=\footnotesize}]
  % becher gauche
  \draw (0,0)--(0,2)--(1.6,2)--(1.6,0)--cycle; \draw[dashed] (0,1.5)--(1.6,1.5);
  \node at (0.8,0.4){\scriptsize$\ch{SO_4^{2-}{+}Zn^{2+}}$};
  \draw[thick] (0.5,1.2)--(0.5,2.6); \node[above] at (0.5,2.6){Zn};
  % becher droit
  \draw (3.4,0)--(3.4,2)--(5,2)--(5,0)--cycle; \draw[dashed] (3.4,1.5)--(5,1.5);
  \node at (4.2,0.4){\scriptsize$\ch{SO_4^{2-}{+}Cu^{2+}}$};
  \draw[thick] (4.5,1.2)--(4.5,2.6); \node[above] at (4.5,2.6){Cu};
  % pont salin
  \draw[thick] (0.5,2.9)to[bend left=45](4.5,2.9); \node[above] at (2.5,3.5){pont salin};
  % voltmetre
  \draw (2.5,3.9) circle (0.28) node{V};
  \draw (0.5,2.6)--(0.5,3.9)--(2.22,3.9); \draw (2.78,3.9)--(4.5,3.9)--(4.5,2.6);
  \node[left] at (0.5,3.4){$-$}; \node[right] at (4.5,3.4){$+$};
  \node at (0.4,-0.35){\scriptsize anode}; \node at (4.6,-0.35){\scriptsize cathode};
\end{tikzpicture}
\end{center}
Demi-pile de gauche (oxydation) : $\ch{Zn \longrightarrow Zn^{2+} + 2\,e^-}$ ; les ions \ch{NO_3^-} migrent vers cette solution. Demi-pile de droite (réduction) : $\ch{Cu^{2+} + 2\,e^- \longrightarrow Cu}$ ; les ions \ch{NH_4^+} migrent vers cette solution (dépôt de cuivre). Le zinc s'oxyde (pôle négatif, anode) ; le cuivre se réduit (pôle positif, cathode). F.é.m. : $E^0(\ch{Cu^{2+}/Cu}) - E^0(\ch{Zn^{2+}/Zn})$.

$\mathbf{a_4 = b_1}$ : le potentiel redox $E^\circ(\text{Ox/Red})$, exprimé en volts, caractérise la force d'un oxydant/réducteur en conditions standard. Plus $E^\circ$ est élevé, plus le pouvoir oxydant de Ox est grand (et plus le pouvoir réducteur de Red est faible).
\begin{center}
\begin{tikzpicture}[>=Stealth,scale=0.9,every node/.style={font=\footnotesize}]
  \draw[->] (0,-1.6)--(0,1.8) node[left]{$E^\circ$};
  \node[rotate=90] at (-1,0){\scriptsize pouvoir oxydant croissant};
  \node[rotate=90] at (4.2,0){\scriptsize pouvoir réducteur croissant};
  \draw (0.2,0.9)--(3,0.9); \node[left] at (0.2,0.9){$E_2^\circ$}; \node at (1.4,1.1){Ox2}; \node at (2.6,1.1){Red2};
  \draw (0.2,-0.6)--(3,-0.6); \node[left] at (0.2,-0.6){$E_1^\circ$}; \node at (1.4,-0.4){Ox1}; \node at (2.6,-0.4){Red1};
  \draw[->] (2.4,0.75)--(1.6,-0.45); \draw[->] (1.6,0.75)--(2.4,-0.45);
  \node at (1.9,-1.3){\scriptsize $\ch{Ox2 + Red1 \longrightarrow Ox1 + Red2}$};
\end{tikzpicture}
\end{center}

\partie{B : application des connaissances}

\question{1. a.} Le nombre de moles de \ch{NH_3} dans $V_1$ (C.N.T.P.) égale celui dans la solution $S_1$ ($V_{S_1} = 2$ L) : $n = \dfrac{V_1}{V_m} = C_{S_1}\times V_{S_1}$, d'où $V_1 = V_m\times C_{S_1}\times V_{S_1}$.

\underline{A.N.} : $V_1 = 22{,}4\times 2{,}08\cdot10^{-4}\times 2 = 9{,}3\cdot10^{-3} \approx 10^{-2}$ L.

\textbf{b.} $\ch{NH_3 + H_2O \rightleftharpoons NH_4^+ + HO^-}$.

\textbf{c.} L'ammoniaque est une base faible (dissociation limitée) ; l'eau en excès est le siège de l'autoprotolyse :
\[ \ch{NH_3 + H_2O \rightleftharpoons NH_4^+ + HO^-} \qquad \ch{2\,H_2O \rightleftharpoons H_3O^+ + HO^-} \]
\underline{Espèces présentes :} ions \ch{H_3O^+}, \ch{HO^-}, \ch{NH_4^+} ; molécules \ch{NH_3}, \ch{H_2O}.

$\bullet$ $[\ch{H_2O}] = \dfrac{\rho_{\text{eau}}}{M_{\ch{H_2O}}} = \dfrac{1000}{18} = 55{,}6\ \text{mol}\cdot\text{L}^{-1}$.

$\bullet$ $[\ch{H_3O^+}] = 10^{-\text{pH}} = 10^{-9{,}7} = 2\cdot10^{-10}\ \text{mol}\cdot\text{L}^{-1}$.

$\bullet$ $[\ch{HO^-}] = \dfrac{K_e}{[\ch{H_3O^+}]} = \dfrac{10^{-14}}{2\cdot10^{-10}} = 5\cdot10^{-5}\ \text{mol}\cdot\text{L}^{-1}$.

$\bullet$ Électroneutralité : $[\ch{H_3O^+}] + [\ch{NH_4^+}] = [\ch{HO^-}]$ ; comme $[\ch{H_3O^+}] \ll [\ch{HO^-}]$, $[\ch{NH_4^+}] = 5\cdot10^{-5}\ \text{mol}\cdot\text{L}^{-1}$.

$\bullet$ Conservation : $C_{S_1} = [\ch{NH_3}] + [\ch{NH_4^+}]$, d'où $[\ch{NH_3}] = 2{,}08\cdot10^{-4} - 5\cdot10^{-5} = 1{,}58\cdot10^{-4}\ \text{mol}\cdot\text{L}^{-1}$.

\textbf{d.} $K_a = \dfrac{[\ch{H_3O^+}][\text{base}]}{[\text{acide}]}$ donne $\text{p}K_a = \text{pH} - \log\dfrac{[\ch{NH_3}]}{[\ch{NH_4^+}]} = 9{,}7 - \log\dfrac{1{,}58\cdot10^{-4}}{5\cdot10^{-5}} = 9{,}2$.

\question{2. a.} Comme pH $=$ p$K_a$, les formes acide et basique ont la même concentration : $S_2$ \textbf{est une solution tampon}. L'ajout d'un peu de \ch{HCl} ($\ch{HCl + H_2O \longrightarrow H_3O^+ + Cl^-}$) est neutralisé par \ch{NH_3} en excès ($\ch{H_3O^+ + NH_3 \longrightarrow NH_4^+ + H_2O}$) : le pH reste quasi constant.

\textbf{b.} À la demi-équivalence, pH $=$ p$K_a$ et $C_{\ch{HCl}}\times V = \dfrac{C_{S_1}\times V'_{S_1}}{2}$ avec $V'_{S_1} = 20$ mL. D'où $V = \dfrac{C_{S_1}\times V'_{S_1}}{2\,C_{\ch{HCl}}} = \dfrac{2{,}08\cdot10^{-4}\times 20\cdot10^{-3}}{2\times 4{,}16\cdot10^{-4}} = 5\cdot10^{-3}$ L $= 5$ mL.

\matiere{PHYSIQUE}{12 points}
\partie{A : vérification des connaissances}

\question{1. Question à réponse courte}
Dans un phénomène d'interférence, l'\textbf{interfrange} est la distance constante qui sépare deux franges brillantes (ou sombres) consécutives.

\question{2. Schéma à compléter}
\begin{center}
\begin{tikzpicture}[>=Stealth,scale=1,every node/.style={font=\footnotesize}]
  \draw[thick] (0,0)--(5.5,2.75); % plan incline
  \draw (0,0)--(5.5,0); \draw (0.9,0) arc[start angle=0,end angle=26.57,radius=0.9];
  % solide (petit rectangle sur le plan)
  \coordinate (M) at (2.6,1.3);
  \draw[rotate around={26.57:(M)}] ($(M)+(-0.5,0)$) rectangle ($(M)+(0.5,0.5)$);
  % poids P vers le bas
  \draw[->,thick] (M) ++(0,0.25) -- ++(0,-1.6) node[below]{$\vv{P}$};
  % composantes
  \draw[->] (M) ++(0,0.25) -- ++(-0.9,0.45) node[left]{$\vv{P_x}$};
  \draw[->] (M) ++(0,0.25) -- ++(0.72,-1.44) node[right]{$\vv{P_y}$};
  \draw[dashed] (M)++(0,0.25)++(-0.9,0.45) -- ++(0.72,-1.44);
  % reaction normale R et tension T
  \draw[->,thick] (M) ++(0,0.55) -- ++(-0.6,1.2) node[above]{$\vv{R}$};
  \draw[->,thick] (M) ++(0.35,0.4) -- ++(0.9,0.45) node[right]{$\vv{T}$};
  % fil vers le haut
  \draw (M)++(0.35,0.4)++(0.9,0.45) -- (5.5,2.75);
  \draw (5.3,2.75)--(5.5,2.95); \draw (5.15,2.6)--(5.35,2.8); \draw (5.45,2.9)--(5.65,3.1);
\end{tikzpicture}
\end{center}

\question{3. Réarrangement}
Dans un circuit \emph{capacitif}, l'impédance du condensateur est supérieure à celle de la bobine. L'impédance d'un circuit $(R, L, C)$ série vaut $Z = \sqrt{R^2 + (X_L - X_C)^2}$ ; le circuit est capacitif si $X_C > X_L$, inductif si $X_L > X_C$.
\begin{center}
\begin{tikzpicture}[>=Stealth,scale=1,every node/.style={font=\footnotesize}]
  \draw (0,0)--(0,1.6) (0,1.6)--(1,1.6) (2,1.6)--(2.6,1.6) (3.6,1.6)--(5,1.6) (5,1.6)--(5,1.1) (5,0.6)--(5,0) (0,0)--(5,0);
  % source
  \draw (0,0.8) circle (0.32); \draw[thick] (-0.15,0.8) sin (0,0.95) cos (0.15,0.8); \node[left] at (-0.5,0.8){$u$};
  % R
  \draw (1,1.45) rectangle (2,1.75); \node[above] at (1.5,1.78){$R$}; \node[below] at (1.5,1.4){$u_R$};
  % L (bobine)
  \draw[decorate,decoration={coil,segment length=4pt,amplitude=4pt}] (2.6,1.6)--(3.6,1.6);
  \node[above] at (3.1,1.78){$L$}; \node[below] at (3.1,1.4){$u_L$};
  % C
  \draw (4.7,1.1)--(5.3,1.1) (4.7,0.9)--(5.3,0.9);
  \node[right] at (5.35,1.0){$C$}; \node at (4.3,1.0){$u_C$};
  \draw[->] (1.6,1.9)--(1.95,1.9); \node[above] at (2.25,1.85){$i$};
\end{tikzpicture}
\end{center}

\partie{B : application des connaissances}

\question{1. a.} De $u(t) = 220\sqrt{2}\cos(100\pi t)$ (V), $U_{\max} = 220\sqrt{2}$ V, d'où $U_{\text{eff}} = \dfrac{U_{\max}}{\sqrt{2}} = 220$ V. Or $U_{\text{eff}} = Z\,I_{\text{eff}}$, donc $Z = \dfrac{U_{\text{eff}}}{I_{\text{eff}}} = \dfrac{220}{2} = 110\ \Omega$.

\textbf{b.} $Z = \sqrt{(R + r)^2 + (L\omega)^2}$ donne $r = \sqrt{Z^2 - (L\omega)^2} - R = \sqrt{110^2 - (0{,}25\times 100\pi)^2} - 50 = 27\ \Omega$.

\textbf{c.} $P_{\text{moy}} = (R + r)I_{\text{eff}}^2 = (50 + 27)\times 2^2 = 308$ W.

\question{2. a.} $I_{\text{eff}}$ est maximale quand $Z = \sqrt{(R+r)^2 + (L\omega - \frac{1}{C\omega})^2}$ est minimale : c'est la \textbf{résonance}, $L\omega = \dfrac{1}{C\omega}$, soit $C = \dfrac{1}{L\omega^2} = \dfrac{1}{0{,}25\times(100\pi)^2} = 4\cdot10^{-5}$ F.

\textbf{b.} À la résonance $Z = R + r = 77\ \Omega$.

\textbf{c.} À la résonance $i$ et $u$ sont en phase : $i(t) = \dfrac{U_{\text{eff}}}{R+r}\sqrt{2}\cos(100\pi t) = \dfrac{220}{77}\sqrt{2}\cos(100\pi t) = 2{,}86\sqrt{2}\cos(100\pi t)$.

\partie{C : résolution d'un problème}
\begin{center}
\begin{tikzpicture}[>=Stealth,scale=1,every node/.style={font=\footnotesize}]
  \draw[->] (-0.3,0)--(8.5,0) node[right]{$x$ (cm)};
  \draw[->] (0,-1.6)--(0,1.9) node[above]{$y$ (mm)};
  \draw[dashed] (0,1.2)--(8,1.2) node[right]{}; \draw[dashed] (0,-1.2)--(8,-1.2);
  \node[left] at (0,1.2){$3$}; \node[left] at (0,-1.2){$-3$};
  \foreach \x in {1,2,3,4}{\draw (2*\x,0.06)--(2*\x,-0.06) node[below]{\scriptsize\x};}
  \draw[thick] plot[domain=0:8,samples=120] (\x,{1.2*sin(deg(pi*\x))});
  \draw[<->,thick] (0,-1.45)--(2,-1.45); \node[below] at (1,-1.45){$\lambda$};
\end{tikzpicture}
\end{center}

\question{1. a.} Par lecture, l'amplitude $a = 3$ mm $= 3\cdot10^{-3}$ m.

\textbf{b.} $\lambda = 2$ cm $= 2\cdot10^{-2}$ m.

\textbf{c.} $c = \lambda N = 2\cdot10^{-2}\times 50 = 1\ \text{m}\cdot\text{s}^{-1}$.

\question{2.} $y_S(t) = a\sin(\omega t + \varphi) = 3\cdot10^{-3}\sin(100\pi t + \varphi)$. Comme $y_S(0) = 0$, $\sin\varphi = 0$ : $\varphi = 0$ ou $\pi$. Le point $S$ partant vers les élongations positives, $y'_S(0) > 0$ : seul $\varphi = 0$ convient. D'où $y_S(t) = 3\cdot10^{-3}\sin(100\pi t)$.

\question{3. a.} $P$ reproduit le mouvement de $S$ avec un retard $t_P = \dfrac{d}{c} = \dfrac{4{,}5\cdot10^{-2}}{1} = 4{,}5\cdot10^{-2}$ s :
\[ y_P(t) = y_S(t - t_P) = 3\cdot10^{-3}\sin(100\pi t - 4{,}5\pi) = 3\cdot10^{-3}\sin\left(100\pi t - \tfrac{\pi}{2}\right). \]

\textbf{b.} $\Delta\varphi = \varphi_P - \varphi_S = -\dfrac{\pi}{2}$ : $P$ est en retard de phase ($\Delta\varphi < 0$) et en quadrature ($|\Delta\varphi| = \frac{\pi}{2}$) par rapport à $S$. Quand $S$ passe par $(0,0)$ en croissant, $P$ passe par son minimum, au point $(4{,}5\,;-3)$.

\end{document}
